\begin{tikzpicture}[font=\small,>=stealth]
\draw[gray!45,rounded corners=3pt] (0,0) rectangle (14.2,10.5);
\node at (7.1,9.85) {焦点在弦内：两段焦半径的倒数相加为定值};

\begin{scope}[shift={(2.4,3.5)},scale=.85]
\draw[->,gray] (-.6,0)--(6,0) node[right] {$x$};
\draw[->,gray] (0,-2.5)--(0,5) node[above] {$y$};
\draw[teal!75!black,thick,samples=90,domain=-2.3:4.5,variable=\v] plot ({\v*\v/4},{\v});

\coordinate (A) at (4,4);\coordinate (B) at (.25,-1);\coordinate (F) at (1,0);
\draw[blue!75!black,very thick] (F)--(A);
\draw[orange!85!black,very thick] (F)--(B);
\foreach \q in {A,B,F}\fill (\q) circle(3pt);
\node[above left=5pt] at (A) {$A$};
\node[left=9pt] at (B) {$B$};\node[below right=6pt] at (F) {$F$};

\node[blue!75!black,above left] at (2.4,2) {$r$};
\node[orange!85!black,below right] at (.5,-.6) {$s$};
\node[below left] at (0,0) {$O$};
\end{scope}
\node[anchor=west,align=left] at (9.1,6.8) {$r=|AF|>0$\\[10pt]$s=|BF|>0$};
\node[anchor=west,align=left] at (9.1,4.7) {图取 $p=2$\\[10pt]$r=5,\ s=5/4$};
\node[anchor=west,align=left] at (9.1,2.6) {$\displaystyle\frac15+\frac45=1$\\[10pt]$\displaystyle\frac2p=1$};
\node at (7.1,.65) {$A,B$ 分居 $F$ 两侧；$|AB|=r+s$，并非把倒数相加当作弦长。};

\end{tikzpicture}
